\chapter{Napovedni model za momenta donosov}
\label{ch:model}

V prejšnjem poglavju smo videli, da moramo za učinkovito sestavo portfelja oceniti momenta donosov, predvsem vektor pričakovanih donosov~$\vec{\mu}$. V praksi za oceno momentov pogosto uporabimo kar drseče zgodovinsko povprečje donosov. Težava tega pristopa pa je, da zgolj naivno ekstrapolira preteklost in se v praksi pogosto izkaže za izjemno slabo napoved. Zato se sodobni kvantitativni pristopi zanašajo na napovedne modele strojnega učenja: analitične postopke, ki se iz velike količine preteklih podatkov samodejno naučijo kompleksnih vzorcev, uporabnih za napovedovanje prihodnjih vrednosti.

\section{Transformer}

Za obdelavo časovnih vrst in relacij med delnicami uporabljamo arhitekturo \emph{transformer}~\cite{vaswani2017}. V preteklosti so se za obdelavo zaporedij uporabljale rekurentne nevronske mreže (kot je LSTM~\cite{hochreiter1997lstm}), ki podatke obdelujejo strogo sekvenčno in zato težko zajamejo dolgoročne odvisnosti, računanja pa ni mogoče dobro vzporediti. Transformer ti omejitvi preseže z mehanizmom \emph{pozornosti} (angl.\ \emph{attention}), ki vsak element vhodnega zaporedja neposredno poveže z vsemi drugimi.

\subsection{Osnovni pojmi}

Ključne pojme transformerja povzemamo po Wangu in sod.~\cite{wang2022}, ki so transformer uporabili za napovedovanje borznih indeksov.

\noindent\textbf{Vložitev in položajno kodiranje} (angl.\ \emph{embedding, positional encoding}): Vhodno zaporedje~$\vec{X}\in\mathbb{R}^{T}$ z \emph{vložitvijo} --- naučeno polno povezano preslikavo --- preslikamo v matriko predstavitev~$\vec{A}\in\mathbb{R}^{T\times d}$. Ker pozornost sama ne loči vrstnega reda elementov, vložitvi prištejemo \emph{položajno kodiranje}~$\vec{PE}\in\mathbb{R}^{T\times d}$, ki časovni položaj zapiše s sinusnimi in kosinusnimi funkcijami različnih frekvenc, kot podaja enačba~\eqref{eq:posenc}:
\begin{equation}
\label{eq:posenc}
  \mathrm{PE}_{(t,2s)}=\sin\!\big(t/10000^{2s/d}\big),\qquad
  \mathrm{PE}_{(t,2s+1)}=\cos\!\big(t/10000^{2s/d}\big).
\end{equation}

\noindent\textbf{Pozornost} (angl.\ \emph{self-attention}): Iz vhodne matrike predstavitev~$\vec{X}\in\mathbb{R}^{n\times d}$ z naučenimi linearnimi preslikavami tvorimo \emph{poizvedbe}~$\vec{Q}$, \emph{ključe}~$\vec{K}$ in \emph{vrednosti}~$\vec{V}\in\mathbb{R}^{n\times d}$; izhod podaja enačba~\eqref{eq:attention}:
\begin{equation}
\label{eq:attention}
  \mathrm{Attention}(\vec{Q},\vec{K},\vec{V}) = \mathrm{softmax}\!\left(\frac{\vec{Q}\vec{K}^{\!\top}}{\sqrt{d}}\right)\vec{V}.
\end{equation}
Skalarni produkt~$\vec{Q}\vec{K}^{\!\top}$ meri podobnost med elementi, funkcija softmax ga pretvori v nenegativne uteži (z vsoto~1), skaliranje z~$\sqrt{d}$ pa prepreči nasičenje softmaxa pri velikih dimenzijah (Slika~\ref{fig:attn}).

\noindent\textbf{Večglava pozornost} (angl.\ \emph{multi-head attention}): Pozornost izvedemo vzporedno v~$h$ glavah, vsako na svoji projekciji poizvedb, ključev in vrednosti, izhode pa zložimo skozi linearni sloj, kot podaja enačba~\eqref{eq:multihead}:
\begin{equation}
\label{eq:multihead}
  \mathrm{MultiHead}(\vec{Q},\vec{K},\vec{V}) = [\vec{H}_1,\dots,\vec{H}_h]\,\vec{W}^{O},\qquad
  \vec{H}_j=\mathrm{Attention}(\vec{Q}\vec{W}_j^{Q},\vec{K}\vec{W}_j^{K},\vec{V}\vec{W}_j^{V}).
\end{equation}
Tako model hkrati zajame odvisnosti iz različnih podprostorov predstavitev.

\noindent\textbf{Kodirnik in dekodirnik} (angl.\ \emph{encoder--decoder}): Izvirni transformer ima kodirnik-dekodirnik zgradbo. \emph{Kodirnik} je sklad~$M$ enakih slojev --- vsak z večglavo pozornostjo in polno povezano podmrežo, ovitima s preskočno povezavo in normalizacijo --- in vhodno zaporedje stisne v predstavitev; \emph{dekodirnik} iz te predstavitve tvori izhodno zaporedje.

\subsection{Delovanje}

V naši izvedbi uporabimo zgradbo \emph{samo s kodirnikom} (angl.\ \emph{encoder-only}): ker ne generiramo izhodnega zaporedja, temveč vsako sredstvo preslikamo v enotno predstavitev~$\vec{z}$, dekodirnika ne potrebujemo. Chen~\cite{chen2025} različne zgradbe (samo kodirnik, samo dekodirnik in kodirnik-dekodirnik) primerja prav na napovedovanju borznih vrednosti, kjer je zgradba samo s kodirnikom preprosta in konkurenčna izbira. Osrednja lastnost pozornosti, ki jo pri tem izkoriščamo, je \emph{permutacijska invariantnost}: ker sama ne loči vrstnega reda enot, je naravna za obdelavo neurejenih množic --- kot je celoten presek delnic ob danem času.

\begin{figure}[htbp]
  \centering
  \includegraphics[width=0.78\linewidth]{transformer_pozornost}
  \caption{Mehanizem pozornosti: poizvedba (rdeča enota) svojo osveženo predstavitev~$\vec{z}$ sestavi kot uteženo vsoto vrednosti vseh enot, kjer so uteži~$\alpha_j$ (debelina puščic) sorazmerne podobnosti med poizvedbo in ključi.}
  \label{fig:attn}
\end{figure}

\newpage

%================================================================
\section{Napovedni transformerji za momenta donosov}
\label{sec:transformerji}

Namesto da bi zgradili lasten napovedni model, v tej nalogi uporabimo \textbf{tri že obstoječe (angl.\ \emph{off-the-shelf}) transformerske modele z javno dostopno kodo}, vsakega z drugačno oblikovno filozofijo, in jih z minimalnimi prilagoditvami vključimo v isti cevovod --- napoved momentov donosov, ki jo nato prevzame optimizator portfelja (Poglavje~\ref{ch:algoritmi}). Takšna zasnova ima tri prednosti: napovedi so \emph{ponovljive} (uporabimo objavljene izvedbe), izognemo se prekomernemu prilagajanju, ki bi ga prinesel bogat lasten model na naši zmerni univerzi sredstev, predvsem pa lahko \emph{pošteno ocenimo}, ali transformerji sploh prispevajo vrednost v primerjavi s klasičnimi ocenami.

Izbrani modeli pokrivajo tri komplementarne osi sodobnega napovedovanja:
\begin{itemize}
  \item \textbf{TACTiS-2}~\cite{ashok2024tactis} --- verjetnostni transformer, ki se nauči \emph{skupne} napovedne porazdelitve nad surovimi donosi vseh sredstev; edini izmed treh napove \emph{oba} momenta ($\vec{\mu}$ \emph{in} $\vec{\Sigma}$) iz iste porazdelitve.
  \item \textbf{MASTER}~\cite{li2024master} --- presečni tržno pogojeni transformer nad bogatim naborom tehničnih značilk; napove razvrstitveni signal $\vec{\mu}$.
  \item \textbf{TFT} (angl.\ \emph{Temporal Fusion Transformer})~\cite{lim2021tft} --- splošni večhorizontni napovedni transformer; napove $\vec{\mu}$.
\end{itemize}
Modela MASTER in TFT napovesta le $\vec{\mu}$; njuno matriko tveganja $\vec{\Sigma}$ ocenimo klasično (vzorčna kovarianca). Vse tri modele nato primerjamo čez več tržnih režimov (Poglavje~\ref{ch:eksperiment}).

\subsection{Osnovni pojmi}

\noindent\textbf{Presečna razvrstitev} (angl.\ \emph{cross-sectional ranking}): Pri izbiri portfelja ni ključna absolutna napoved donosa vsakega posameznega sredstva, temveč njihova relativna ureditev --- katera sredstva bodo v prihodnjem obdobju \emph{relativno} boljša (ali slabša) od drugih v naboru. Ker so donosi sredstev izrazito soodvisni prek skupnih makroekonomskih, tržnih in sektorskih dejavnikov, napredni modeli obravnavajo celoten presek sredstev \emph{hkrati} in med njimi neposredno modelirajo odvisnosti. Presečna razvrstitev tako preusmeri fokus iz napovedovanja točne številke na relativno hierarhijo, kar neposredno služi optimizatorju: za uspešno alokacijo kapitala ni pomembno, koliko je trg zrasel v absolutnem smislu, temveč katera sredstva nosijo največji relativni presežni donos (angl.\ \emph{alpha}).\\

\noindent\textbf{Informacijski koeficient} (IK; angl.\ \emph{information coefficient})~\cite{grinold2000}: Osrednja mera kakovosti napovedi, definirana kot presečna rang-korelacija med napovedanim signalom modela in dejansko uresničenim prihodnjim donosom čez celoten presek sredstev. IK neposredno meri prav to, kar za optimizacijo portfelja šteje: kako zanesljivo model loči zmagovalna sredstva od tistih, ki bodo prinesla izgubo.\\

\noindent\textbf{Verjetnostna napoved in kopula} (angl.\ \emph{probabilistic forecast, copula}): Namesto ene točkovne ocene lahko model napove celotno \emph{porazdelitev} prihodnjih donosov, iz katere lahko vzorčimo. Za več sredstev hkrati je poleg robnih porazdelitev posameznih donosov treba modelirati še njihovo \emph{soodvisnost}; ta struktura odvisnosti se matematično zajame s \emph{kopulo}. Iz vzorcev skupne porazdelitve lahko neposredno ocenimo oba momenta: $\vec{\mu}$ kot povprečje in $\vec{\Sigma}$ kot kovarianco vzorcev.\\

\noindent\textbf{Tržno pogojena vrata} (angl.\ \emph{market-guided gating}): Vektor dinamičnih uteži, ki ga model sproti izračuna iz \emph{tržnega konteksta} (stanje širokih indeksov, volatilnost, obseg trgovanja). Deluje kot inteligentni filter, ki glede na trenutni tržni režim poudari informativne dimenzije vhodnih predstavitev in zaduši prešumne.\\

\newpage

\subsection{TACTiS-2: transformer s pozornostno kopulo}
\label{sec:tactis}
Model \textbf{TACTiS-2} (angl.\ \emph{Transformer-Attentional Copulas for Time Series}) Ashoka in sod.~\cite{ashok2024tactis} je verjetnostni napovedni transformer, ki se nauči \emph{skupne} večrazsežne porazdelitve prihodnjih donosov vseh sredstev hkrati. Robne porazdelitve posameznih donosov modelira z normalizacijskim tokom, njihovo soodvisnost pa s \emph{pozornostno kopulo} --- kopulo, katere parametre generira mehanizem pozornosti. Model se uči v dveh fazah: najprej se naučijo robne porazdelitve, nato še kopula odvisnosti.

Momenta donosov iz modela pridobimo neposredno z vzorčenjem. Ob izbranem času $t$ modelu podamo okno zadnjih $\sim\!45$ dni surovih dnevnih donosov vseh sredstev; iz naučene skupne porazdelitve nato vzorčimo $S\!\approx\!400$ skupnih poti donosov za naslednjih $H$ dni. Vsako pot seštejemo v kumulativni $H$-dnevni donos, s čimer dobimo matriko vzorcev $\vec{C}\in\mathbb{R}^{N\times S}$. Momenta ocenimo kot vzorčno statistiko, kot podaja enačba~\eqref{eq:tactismoments}:
\begin{equation}
\label{eq:tactismoments}
  \vec{\mu} = \tfrac{1}{S}\textstyle\sum_{s} \vec{C}_{:,s}, \qquad
  \vec{\Sigma} = \operatorname{Cov}\!\big(\vec{C}\big).
\end{equation}
TACTiS-2 je torej \emph{edini} od naših modelov, ki poda oba momenta iz iste napovedane porazdelitve --- $\vec{\mu}$ z resnično magnitudo in $\vec{\Sigma}$, ki je z njo usklajena --- ne da bi karkoli prevzel iz zgodovinske ocene. Uporabimo ga prek javno dostopnega paketa \texttt{tactis}, kar zahteva le minimalne prilagoditve.

\subsection{MASTER: tržno pogojeni presečni transformer}
\label{sec:master}
Model \textbf{MASTER} (angl.\ \emph{Market-Guided Stock Transformer}) Lija in sod.~\cite{li2024master} je presečni transformer, zasnovan prav za napovedovanje razvrstitve donosov na celotnem preseku delnic. Uporabimo \emph{izvirno} izvedbo iz javnega repozitorija; njene vhodne značilke zgradimo sami iz cen in obsega trgovanja (OHLCV), tako da model teče na naši univerzi brez posebne podatkovne infrastrukture.

Vhod ob času $t$ za vsako sredstvo sestavlja okno zadnjih $L=8$ dni s $158$ tehničnimi značilkami (nabor \emph{Alpha158}: drseča povprečja, momentum, volatilnost, količinska razmerja ter korelacije cene z obsegom čez več časovnih oken) in $63$ tržnih značilk, ki napajajo tržno pogojena vrata (izpeljemo jih iz treh širokih indeksov). Model surove značilke najprej pretvori v predstavitev in jih modulira s tržnimi vrati, nato izvede \emph{časovno} pozornost (vzorci znotraj posamezne delnice), \emph{med-sredstveno} pozornost (relacije čez celoten presek delnic) in časovno agregacijo, na koncu pa iz predstavitve $\vec{z}_i$ z linearno preslikavo napove skalarni razvrstitveni signal $s_i$. Model se uči z minimizacijo srednje kvadratne napake na presečno standardiziranem prihodnjem donosu, kar v praksi ustreza razvrstitvenemu cilju. MASTER napove le $\vec{\mu}$ (razvrstitev); matriko tveganja $\vec{\Sigma}$ zanj ocenimo klasično.

\subsection{TFT: časovni fuzijski transformer}
\label{sec:tft}
Model \textbf{TFT} (angl.\ \emph{Temporal Fusion Transformer}) Lima in sod.~\cite{lim2021tft} je splošni večhorizontni napovedni transformer, ki ni zasnovan posebej za finance --- prav zato je pomemben preizkus, ali \emph{splošni} transformer sploh prispeva vrednost. Združuje mreže za \emph{izbor spremenljivk} (samodejno tehtajo pomen vhodnih značilk), \emph{vrata} za nadzor pretoka informacij, interpretabilno večglavo pozornost ter \emph{kvantilni} izhod (napove razpon, ne le točkovne vrednosti). Uporabimo ga prek javnega paketa \texttt{pytorch-forecasting}.

Vsakemu sredstvu podamo zaporedje dnevnih donosov z nekaj preprostimi kazalniki (momentum, volatilnost) in statično vložitev identitete delnice. Model napove \emph{naslednjih $H$ dnevnih donosov}; napovedani $H$-dnevni pričakovani donos posameznega sredstva dobimo kot vsoto teh napovedi (enaka definicija kot pri TACTiS-2). TFT napove le $\vec{\mu}$; $\vec{\Sigma}$ zanj ocenimo klasično.

\subsection{Vključitev v cevovod}
\label{sec:vkljucitev}
Modeli se razlikujejo v tem, kaj napovejo, zato jih v cevovod vključimo na dva načina. Za \textbf{TACTiS-2} $\vec{\mu}$ in $\vec{\Sigma}$ vzamemo neposredno iz napovedane porazdelitve (enačba~\eqref{eq:tactismoments}); nič ni prevzeto iz zgodovine. Za \textbf{MASTER} in \textbf{TFT} pa velja pomembna previdnost: njun signal je presečno \emph{normiran} in ne poda absolutne ravni donosa. Zato ga presečno standardiziramo in afino reskaliramo na povprečje ter razpršenost istočasne zgodovinske ocene~$\vec{\mu}$, kot povzema enačba~\eqref{eq:rescale}:
\begin{equation}
\label{eq:rescale}
  \vec{\mu}_{\text{model}} = \operatorname{zscore}(\vec{s})\cdot \sigma_{\vec{\mu}_{\text{hist}}} + \bar{\mu}_{\text{hist}}.
\end{equation}
S tem sta absolutna raven in razpon napovedi (in torej razmerje med donosom in tveganjem, ki ga vidi optimizator) enaka kot pri zgodovinskem scenariju; med scenariji se spreminja le presečni vrstni red sredstev, torej \emph{katera} sredstva so relativno privlačnejša. Prispevek teh dveh modelov je torej boljša razvrstitev, ne absolutna napoved donosa. Njuno matriko tveganja $\vec{\Sigma}$ ocenimo z vzorčno kovarianco. Za pošteno primerjavo TACTiS in TFT uporabljata \emph{enako} dolžino vhodnega okna (45 dni) kot drseči zgodovinski baseline; MASTER ohrani svoje izvirno okno 8 dni, saj njegove značilke Alpha158 v vsak korak že vgradijo daljšo zgodovino (drseče statistike do 60 dni).

V vseh primerih optimizator napovedi ne uporabi samostojno: pod kardinalno omejitvijo (največ $K$ sredstev) razvrstitev~$\vec{\mu}$ tehta proti oceni tveganja~$\vec{\Sigma}$ (Poglavje~\ref{ch:algoritmi}). Pri parametru tveganja~$P\!\to\!1$ bi sledil zgolj razvrstitvi (izbral najbolje uvrščena sredstva), sicer pa jo umeri z manjšanjem tveganja portfelja. Ključne lastnosti treh modelov povzema Tabela~\ref{tab:hp-model}.

\begin{table}[htbp]
  \centering
  \caption{Primerjava treh napovednih transformerjev.}
  \label{tab:hp-model}
  \small
  \begin{tabular}{lccc}
    \toprule
     & \textbf{TACTiS-2} & \textbf{MASTER} & \textbf{TFT}\\
    \midrule
    vir / paket        & \texttt{tactis} & izvirni repozitorij & \texttt{pytorch-forecasting}\\
    referenca          & \cite{ashok2024tactis} & \cite{li2024master} & \cite{lim2021tft}\\
    vhodne značilke    & surovi donosi & Alpha158 + tržne & donosi + kazalniki\\
    dolžina okna       & 45 dni & 8 dni & 45 dni\\
    izhod              & porazdelitev & razvrstitev & kvantili donosov\\
    napove $\vec{\mu}$ & da (vzorčno povpr.) & da (razvrstitev) & da (vsota donosov)\\
    napove $\vec{\Sigma}$ & \textbf{da} (vzorčna kov.) & ne & ne\\
    vključitev $\vec{\mu}$ & neposredno & reskaliran na zgod. & reskaliran na zgod.\\
    napovedni horizont & $H$ dni & $H$ dni & $H$ dni\\
    \bottomrule
  \end{tabular}
\end{table}
